\documentclass[11pt]{article}
\usepackage[margin=0.8in]{geometry}
\usepackage{amsmath,microtype}
\newcommand{\guideheading}[1]{\par\medskip\noindent{\large\bfseries #1}\par\smallskip}
\begin{document}
\begin{center}
{\Large Guide: Which Late Bridges Cannot Work}\par\medskip
Collatz proof project\quad---\quad September 5, 2026
\end{center}
\guideheading{The search being studied}
Our transfer search follows a known-convergent source and occasionally uses
a smaller transfer instance to introduce an auxiliary source. At the level
of base values, the bridge replaces $x$ by $9x+2$. The aim is to make a
source meet the target uniformly across a whole arithmetic progression.

\guideheading{The stronger result}
Lean now proves a size bound for every finite path combining ordinary
forward Collatz steps and these bridges, in any order. The proof does not
limit the number of bridges or require them to occur at a particular value.
It even allows more bridges than the actual induction rules permit, so
its impossibility consequence applies to the restricted search too.

\guideheading{Why the bound is simple}
An even step halves the value. An odd shortcut step increases it by at most
a factor of two. A bridge increases it by at most a factor of sixteen.
Multiplying those bounds gives a constraint on the endpoint after any
combination of operations. Lean checks the multiplication argument for
arbitrary paths.

\guideheading{What happens at the trivial cycle}
Once a base value is 1 or 2, we can combine the clock, its cycle phase, and
the symbolic slope exponent into an integer quantity called its charge.
Ordinary steps around the cycle preserve this charge.
The new bound proves that a source cannot lower its charge by leaving the
cycle through bridges and eventually returning, even if it inserts more
bridges during the excursion.

If the source charge is already above the target charge, it cannot match
the target later using these forward operations. Increasing the future
bridge budget will not fix that mismatch.

\guideheading{How this improves the search}
The program now discards such states immediately. A complete comparison
for parameters 0 through 1000 found the same certificate outcomes and
first successful depths with this pruning switched on or off. All eight
affine-search tests passed. The pruning rule rests on a general Lean
inequality, not on the finite comparison alone.

\guideheading{What remains possible}
The result does not exclude auxiliary paths introduced earlier, inverse
orbit steps, or different proof rules. It also does not say that the
excluded parameters fail to converge. They are exclusions from this
particular search strategy after it reaches a particular state.

The full Collatz goal still needs a proof covering every parameter.
These findings tell us where adding more of the same forward bridges
cannot help; they do not supply the missing global argument.

\guideheading{Where the proof lives}
Build \texttt{Collatz.BridgeGrowth} from \texttt{lean/}. Its three public
results are included in the theorem axiom audit. The technical paper gives
the exact inequality and its application to symbolic endpoint matching.
\end{document}
